\subsection{Expected Results and Required Instruments}

After the quantity to be measured has been established, the expected range - or at least the order of magnitude - of the quantity should be estimated. This is important for selecting appropriate instruments and settings. In many cases, this process is exploratory, involving small probing measurements to gain insight into the expected behavior. For some measurements, this can be straightforward; for example, measuring a DC voltage with a general-purpose voltmeter is usually simple, as these instruments are well-protected, robust, and reliable.

RF measurements, such as those involving spectrum analyzers or vector network analyzers (VNAs), are considerably more complex. These instruments are expensive and sensitive, and misuse can easily damage them. Understanding the expected magnitude of the signal is therefore critical. Spectrum analyzer measurements, for example, are strongly affected by settings such as resolution bandwidth, attenuator setting, pre-amplifier configuration, and reference level. Incorrect settings can alter the internal measurement chain, causing inaccurate readings, increased intermodulation or distortion, spurious harmonics, or a raised noise floor.

By estimating the expected results in advance, the experimenter ensures that instruments are configured appropriately, measurement quality is optimized, and equipment is protected from potential damage. Admittedly, this can be a tedious task at times. Once the approximate magnitude of the quantity is known, the experimenter should consult the instrument manual to verify proper settings or seek guidance from an experienced technician. However, it is important to do preliminary research before asking for help, in order to use the expert’s time efficiently.

Another important reason to approximate the range of the quantity to be measured is that it allows for a preliminary estimation of measurement uncertainty. This early assessment helps determine whether the chosen method and instruments are suitable for the task at hand, or whether an alternative approach is required. By considering the potential uncertainty upfront, the experimenter can avoid investing time in measurements that are unlikely to yield meaningful or reliable results.

As an extension of this, the expected magnitude and required quality of the measurement should guide instrument selection. For example, highly accurate DC voltage measurements may exceed the capabilities of a standard multimeter. Similarly, for AC voltages, if the frequency is in the hundreds of kilohertz or higher, most general-purpose multimeters will not provide reliable readings. While an oscilloscope is very accurate in the time domain, it is not necessarily precise in amplitude, particularly for small signals.

By considering both the expected magnitude and the required measurement fidelity, the experimenter can choose instruments that are suitable for the task and avoid misinterpretation of results caused by inappropriate equipment. Once again, it is strongly recommended to consult an experienced technician before using any non-general-purpose instruments. These devices are often expensive and sensitive, and improper use can result in damage or inaccurate measurements.